% !PRINTED 407
\weberpage{407}
If $\tau', \tau''$ are any two of these roots, then from (3) it follows that:
\[
\begin{aligned}
\varphi_h\,(\tau')\,\psi_h\,(\tau) - \varphi_h\,(\tau)\,\psi_h\,(\tau') &= 0\\
\varphi_h\,(\tau'')\,\psi_h\,(\tau) - \varphi_h\,(\tau)\,\psi_h\,(\tau'') &= 0,
\end{aligned}
\]
and from this, since $\psi_h\,(\tau)$ and $\varphi_h\,(\tau)$ cannot vanish
at the same time,
\[
\varphi_h\,(\tau')\,\psi_h\,(\tau'') - \varphi_h\,(\tau'')\,\psi_h\,(\tau') = 0,
\]
i.\,e., if $\tau', \tau''$ are any two of the quantities (7), we have
\[
\Phi_h\,(\tau', \tau'') = 0.
\]

From this it follows that the $\nu$ functions $\Phi_h\,(t,\tau')$, apart from a
factor independent of $t$, have the same greatest common divisor as the
functions $\Phi_h\,(t,\tau)$. The same holds for the functions
$\Phi_h\,(t,\tau'')$ and so on, or the equations
\[
R\,(t,\tau) = 0, \quad R\,(t,\tau') = 0, \quad R\,(t,\tau'') = 0, \ldots
\]
all have the same roots. From this it follows, by (5):
\[
\vartheta = \dfrac{R\,(a,\tau)}{R\,(b,\tau)} = \dfrac{R\,(a,\tau')}{R\,(b,\tau')} = \dfrac{R\,(a,\tau'')}{R\,(b,\tau'')} = \cdots,
\]
and consequently, by (6), the $\mu$ values
\[
\xi = \tau, \tau', \tau'', \ldots
\]
are the roots of the equation $X = 0$. On the other hand, from
$\Phi_h\,(\tau',\tau) = 0$, $\Phi_h\,(\tau'',\tau) = 0, \ldots$ it follows that
\[
v_h = \dfrac{\varphi_h\,(\tau)}{\psi_h\,(\tau)} = \dfrac{\varphi_h\,(\tau')}{\psi_h\,(\tau')} = \dfrac{\varphi_h\,(\tau'')}{\psi_h\,(\tau'')} \cdots = \dfrac{1}{\mu}\,\Sigma\,\dfrac{\varphi_h\,(\tau)}{\psi_h\,(\tau)},
\]
and hence $v_h$, as a symmetric function of the roots of (6), can be
represented rationally in terms of the coefficients of that equation,
i.\,e. rationally in terms of $\vartheta$. If we again interchange $t$ with
$\tau$, we obtain by (1) and (2) expressions of the form

\weq{10}{\begin{aligned}
u &= f\,(\vartheta), \quad u_1 = f_1\,(\vartheta), \ldots, u_{\nu-1} = f_{\nu-1}(\vartheta)\\
\vartheta &= \chi\,(u, u_1, \ldots, u_{\nu-1}),
\end{aligned}}

in which $f, f_1, \ldots, f_{\nu-1}, \chi$ denote rational functions.

\begin{center}
\S.\ 108.\\
\gesperrt{Resolvents with one parameter.}
\end{center}

We now return to our initial assumptions (\S.\,106), and denote by
$u, u_1, \ldots, u_{\nu-1}$ a system
%CONT>>
